% !TeX root = ../教材题库_全部条目_转写工作稿.tex
% Source: XJ-B1-C05-EX-P193-Y0513
% 第5章 / 5.4 函数 y=A sin(ωx+φ) 的图象与性质 / 教材 p.193、194 / PDF 第200、201页
% OCR未校对：公式与文字均可能有误；未完成数学TeX转写；禁止直接发给学生
% 原题及相邻原解截图（核对用）：../题目原图/XJ-B1-C05-EX-P193-Y0513-01.webp, ../题目原图/XJ-B1-C05-EX-P193-Y0513-02.webp
% 本文件只在转写完成后作为题干使用；OCR 条目尚不是可靠题干。
\begin{bankitem}{5.4 例题 例题5}{第5章 / 5.4 函数 y=A sin(\ensuremath{\omega}x+\ensuremath{\varphi}) 的图象与性质 / 教材 p.193、194 / PDF 第200、201页}{XJ-B1-C05-EX-P193-Y0513}{OCR未校对：公式与文字均可能有误；未完成数学TeX转写；禁止直接发给学生}
\begin{ocrdraft}
[教材p.193 / PDF第200页]
\par
00画出函数3一2（2，十的简图，并求出这个函数的周期和值域．
\par
（方法一）函数y=2sin2，十了）的周期T=
\par
解
\par
，我们先用“五点法”作
\par
2
\par
出它在一个周期内的图象．
\par
令2十
\par
=0，得x=
\par
玉作为第一点的横坐标，依次递加一个周
\par
8
\par
8
\par
期的+，即匹
\par
，就可以得到其余四个点的横坐标．列表如下：
\par
3丌
\par
5丌
\par
8
\par
8
\par
8
\par
8
\par
8
\par
3丌
\par
2工+
\par
0
\par
2
\par
2
\par
3，一2sin2+
\par
0
\par
2
\par
0
\par
0
\par
作出函数在〔一
\par
丌7丌
\par
一]上的简图
\par
，并逐次向左和向右平移丌个单位长度，就可以得
\par
8
\par
函数3一2（2，十了）的值域为[一2，2]．
\par
到这个函数的图象，如图5．4一8．
\par
y=2sin(2x+a)
\par
2
\par
一匹0
\par
3是
\par
8
\par
8
\par
8
\par
图5．4一8
\par
（方法二）先作出函数3-sinx的图象，将正弦曲线向左平移《个单位长度，
\par
得
\par
到函数！y一（，十《）的图象．
\par
[教材p.194 / PDF第201页]
\par
然后将3一“，（，十）的图象上每一点的纵坐标不变、横坐标缩短为原来的
\par
1
\par
得到函数3一"（2，十了）的图象．
\par
再将函数！y一"，（2，十的图象上每一点的横坐标不变、纵坐标扩大为原来的2
\par
倍，就得到函数3，
\par
=2sin2，十《）的图象，如图5．4一9．
\par
y一2sin（2十）
\par
2
\par
y=sin(2x+\ocrglyph{00E5})
\par
还有其他的变换思路
\par
7訌
\par
42是
\par
吗？试一试．
\par
3訌
\par
8
\par
8
\par
y=sln
\par
Y=sin(X+—
\par
一2
\par
图5．4一9
\par
函数y=2sin2，十《）的周期为丌
\par
，值域为[一2，2]．
\par
现在，我们回过头来重新审视本节开头提到的简谐振动的图象（图5．4一2、图
\par
5，4一3).这些图象所对应的函数解析式为3，一，Asin（-十@的形式，而式子中的参
\par
数在物理学中有明确的意义．例如，简谐振动中，A（A>0）表示这个振动物体偏离
\par
平衡位置的最大距离，称为振幅．
\par
若工表示时间()e[0，十亠）），则这个简谐振动的周期是T=
\par
，其中>0，
\par
而丆一
\par
里表示单位时间内往复振动的次数，称为频率-十称为相位．
\par
工=0
\par
T2丌
\par
时的相位称为初相．
\end{ocrdraft}
\end{bankitem}
